\chapter{Septoplasty}

\section*{What Is This Surgery?}
Septoplasty is a common surgical procedure performed by otolaryngologists to correct a deviated nasal septum. The nasal septum is the wall of cartilage and bone that divides the nasal cavity into two halves. When this septum is significantly displaced or bent, it can obstruct nasal airflow, leading to chronic nasal obstruction, difficulty breathing, recurrent sinusitis, and epistaxis. The primary goal of septoplasty is to straighten the cartilage and bone of the septum, thereby widening the nasal passages and restoring proper respiratory function. This procedure is typically performed entirely through incisions made inside the nose (endonasally), ensuring no external scars and aiming to significantly improve a patient's quality of life by enhancing nasal patency.

\section*{Alternative Names}
\begin{itemize}
  \item Nasal septal surgery
  \item Septal reconstruction
  \item Submucous resection of the nasal septum (SMR)
  \item Septal repair
  \item Functional septoplasty
  \item Septorhinoplasty (when combined with external nasal reshaping)
\end{itemize}

\section*{Relevant Surgical Anatomy}
The nasal septum is a complex osteocartilaginous structure forming the medial wall of the nasal cavity. Its anterior and superior portion is composed of the quadrangular cartilage, which articulates with the perpendicular plate of the ethmoid bone superiorly and the vomer and maxillary crest inferiorly and posteriorly. The perpendicular plate of the ethmoid bone forms the superior-posterior part of the bony septum, while the vomer forms the inferior-posterior part, resting on the maxillary and palatine bones. The entire septum is covered by a delicate mucoperichondrial (over cartilage) and mucoperiosteal (over bone) membrane, which is crucial for its blood supply and integrity.

The rich vascular supply to the nasal septum is primarily derived from branches of both the internal and external carotid arteries. A critical anastomotic area, Kiesselbach's plexus (Little's area), is located in the anterior inferior part of the septum and is formed by contributions from the anterior ethmoidal artery (from the ophthalmic artery), posterior ethmoidal artery (from the ophthalmic artery), sphenopalatine artery (from the maxillary artery), greater palatine artery (from the maxillary artery), and superior labial artery (from the facial artery). Venous drainage largely parallels the arterial supply. Sensory innervation is provided by branches of the trigeminal nerve, specifically the anterior ethmoidal nerve (from the ophthalmic division, V1) and the nasopalatine nerve (from the maxillary division, V2). Lymphatic drainage typically flows to the submandibular and deep cervical lymph nodes. Key surgical landmarks include the anterior septal angle, the dorsal and caudal septal struts, the maxillary crest, and the perpendicular plate of the ethmoid. Preservation of the mucoperichondrial flaps and the L-strut (dorsal and caudal cartilage) is paramount to prevent complications such as septal perforation or saddle nose deformity.

\section*{Surgical Principle \& Rationale}
The fundamental surgical principle of septoplasty involves the meticulous elevation of the mucoperichondrial and mucoperiosteal flaps from the underlying deviated septal cartilage and bone. This creates a subperichondrial/subperiosteal plane, providing direct access to the structural deformities while preserving the vital mucosal lining. Once exposed, the surgeon can then reshape, reposition, or selectively resect the obstructive portions of the septum. The primary rationale is to alleviate mechanical obstruction to nasal airflow by creating a straight, midline septal partition, thereby improving nasal breathing and associated symptoms.

The technical goals of septoplasty are multifaceted. Firstly, it aims to straighten the cartilaginous and bony septum to establish a patent airway on both sides of the nasal cavity. Secondly, it is crucial to preserve the integrity of the septal mucosa to prevent postoperative septal perforation, which can lead to crusting, whistling, and recurrent epistaxis. Thirdly, maintaining adequate septal support, particularly the dorsal and caudal cartilaginous struts (the "L-strut"), is essential to prevent external nasal deformities such as a saddle nose or changes in tip projection. Finally, the procedure often addresses concomitant issues like septal spurs impinging on turbinates or, if performed in conjunction with other procedures, facilitates access for conditions such as turbinate hypertrophy or chronic rhinosinusitis. Techniques employed range from scoring the cartilage to release its intrinsic memory, morselization, limited resection of severely deviated segments, to more complex extracorporeal septoplasty where the septum is removed, reshaped on a back table, and then reimplanted. Following structural correction, the mucosal flaps are carefully re-approximated and stabilized to promote healing and prevent septal hematoma.

\section*{History \& Surgical Evolution}
The surgical correction of septal deviation has a rich history, evolving significantly from rudimentary and often destructive approaches to sophisticated, reconstructive techniques. Early attempts in the 19th century were frequently aggressive, involving extensive resection of the septum, which, while addressing obstruction, often led to severe complications such as septal perforation and the disfiguring saddle nose deformity due to loss of structural support.

A major turning point occurred in the late 19th and early 20th centuries with the development of the submucous resection (SMR) technique. In 1882, Ingals described a technique for septal resection, but it was Arthur E. Freer in 1902 and Gustav Killian in 1904 who independently refined the SMR procedure. Their innovation emphasized the preservation of the overlying mucoperichondrial flaps, which significantly reduced the incidence of perforation and improved healing. Killian's submucous resection became the standard for many decades, involving the elevation of mucosal flaps and removal of the deviated cartilage and bone. While effective for severe deviations, this technique could still compromise septal support if too much cartilage was resected. Modern septoplasty techniques, largely developed in the mid-20th century and refined thereafter, prioritize cartilage preservation and reconstruction over extensive resection. Surgeons like Maurice Cottle advocated for conservative septal surgery, focusing on repositioning and reshaping rather than extensive removal, thereby minimizing the risk of external nasal deformities and preserving the structural integrity of the nose. The advent of endoscopic techniques in the late 20th century further enhanced precision and visualization, allowing for more targeted and less invasive corrections, making septoplasty a highly refined and successful procedure today.

\section*{Clinical Indications}
Septoplasty is primarily indicated to alleviate symptoms caused by a deviated nasal septum that significantly impairs nasal function and quality of life, and which has not responded to conservative medical management.

\begin{itemize}
  \item Chronic nasal obstruction that is unresponsive to medical management (e.g., nasal steroids, antihistamines, decongestants).
  \item Difficulty breathing through the nose, especially during sleep, physical exertion, or in specific positions.
  \item Recurrent or chronic rhinosinusitis due to impaired sinus drainage caused by septal deviation, often confirmed by CT imaging.
  \item Frequent nosebleeds (epistaxis) resulting from turbulent airflow over a deviated septum, leading to mucosal drying, crusting, and irritation.
  \item Snoring or sleep-disordered breathing, where nasal obstruction is a significant contributing factor to increased airway resistance.
  \item Facial pain or headaches attributed to septal spurs impinging on nasal turbinates or other sensitive intranasal structures.
  \item As an access procedure to facilitate other nasal or sinus surgeries, such as endoscopic sinus surgery or turbinate reduction.
  \item Aesthetic concerns when the septal deviation contributes to an external nasal deformity, often addressed in conjunction with rhinoplasty (septorhinoplasty).
\end{itemize}

\section*{Clinical Positioning}
Septoplasty is considered a primary surgical treatment for symptomatic nasal septal deviation. It is typically indicated when conservative medical therapies, such as nasal saline rinses, topical nasal corticosteroids, and oral decongestants, have failed to provide adequate relief from nasal obstruction or other related symptoms. In most cases, a significant, symptomatic septal deviation represents a mechanical obstruction that cannot be effectively managed by non-surgical means. Therefore, once a clear, symptomatic deviation is identified and medical options are exhausted or deemed inappropriate, septoplasty becomes the first-line surgical intervention. It may also be performed as a primary procedure in conjunction with other surgeries, such as turbinate reduction or endoscopic sinus surgery, to optimize overall nasal and sinus function. In rare instances, it might be considered a secondary procedure if a previous septoplasty was unsuccessful, if new septal deviation develops following trauma, or if there is a recurrence of deviation due to cartilage memory or incomplete initial correction.

\section*{Surgical Technique \& Procedure}
Septoplasty is typically performed under general anesthesia, though local anesthesia with sedation can be used in select cases, particularly for minor anterior deviations. The patient is positioned supine with the head slightly elevated, often in a mild reverse Trendelenburg position, to optimize surgical field visualization and minimize venous congestion.

The key steps of the procedure are as follows:
\begin{enumerate}
  \item \textbf{Anesthesia and Vasoconstriction}: General anesthesia is induced. The nasal mucosa is then infiltrated with a local anesthetic containing a vasoconstrictor (e.g., 1\% lidocaine with 1:100,000 epinephrine) to achieve profound local anesthesia, minimize intraoperative bleeding, and provide postoperative pain control. Topical vasoconstrictors may also be applied.
  \item \textbf{Incision}: A small incision, typically a hemitransfixion incision or a Killian's incision, is made just inside the nostril on one side of the septum, usually on the more deviated side. This incision is strategically placed to be hidden and does not result in external scarring.
  \item \textbf{Mucoperichondrial Flap Elevation}: The surgeon carefully elevates the mucoperichondrial and mucoperiosteal flaps from the underlying septal cartilage and bone using specialized elevators. This creates a tunnel-like space, exposing the deviated structures while preserving the vital mucosal lining. Elevation is performed on one side initially, then often on the contralateral side through the same incision or a separate contralateral incision, creating a bipedicled flap.
  \item \textbf{Identification and Correction of Deviation}: The specific areas of deviation, such as C-shaped or S-shaped bends, spurs, ridges, or caudal dislocations in the cartilage and bone, are identified. The surgeon then uses specialized instruments (e.g., septal scissors, osteotomes, rongeurs) to straighten, reshape, or selectively remove the obstructive portions.
  \item \textbf{Cartilage and Bone Reshaping/Resection}: Deviated cartilage may be scored to release its intrinsic memory, morselized, or partially resected to allow it to straighten. Bony spurs or ridges, particularly along the maxillary crest or perpendicular plate, are typically removed. Extreme care is taken to leave a dorsal and caudal strut of cartilage (at least 1 cm wide, often referred to as the "L-strut") to maintain nasal support and prevent a saddle nose deformity or tip ptosis. For severe, complex deviations, an extracorporeal septoplasty may be performed where the entire septum is removed, reshaped ex vivo, and then reimplanted.
  \item \textbf{Flap Re-approximation and Stabilization}: Once the septum is straightened and repositioned, the mucosal flaps are carefully repositioned against the newly aligned septum. It is crucial to ensure smooth re-approximation without tension or dead space.
  \item \textbf{Closure and Support}: The initial incision is closed with fine absorbable sutures. To stabilize the septum, prevent septal hematoma, and promote proper healing, internal splints (thin silicone sheets) may be placed on either side of the septum and secured with through-and-through sutures. Alternatively, quilting sutures may be used to re-approximate the mucosal flaps directly to the septum without splints. Nasal packing, if used, is typically soft and dissolvable or removed within 24-48 hours.
\end{enumerate}
The entire procedure typically takes 30 to 90 minutes, depending on the complexity of the deviation and whether concomitant procedures are performed.

\section*{Preoperative Workup \& Preparation}
Thorough preoperative preparation is crucial for a successful septoplasty, minimizing risks, and optimizing patient outcomes. This process ensures the patient is medically fit for surgery and understands the procedure.

Key aspects of preoperative workup and preparation include:
\begin{itemize}
  \item \textbf{Medical Evaluation}: A comprehensive medical history and physical examination are performed, including a detailed assessment of the nasal passages via anterior rhinoscopy and often nasal endoscopy to visualize the extent and nature of the septal deviation.
  \item \textbf{Imaging}: While not always necessary for straightforward cases, a computed tomography (CT) scan of the paranasal sinuses may be ordered if there is suspicion of concomitant sinus disease, for complex anatomical assessment, or to rule out other causes of nasal obstruction.
  \item \textbf{Laboratory Tests}: Routine preoperative blood tests, such as a complete blood count (CBC), basic metabolic panel, and coagulation studies (PT/INR, PTT), are typically performed to assess overall health, kidney function, and bleeding risk.
  \item \textbf{Medication Review}: Patients are advised to discontinue blood-thinning medications (e.g., aspirin, NSAIDs, warfarin, novel oral anticoagulants) for a specified period (typically 7-14 days) before surgery to significantly reduce the risk of intraoperative and postoperative bleeding. Herbal supplements with anticoagulant properties (e.g., ginkgo biloba, garlic, ginseng) should also be stopped.
  \item \textbf{Anesthesia Clearance}: Patients undergoing general anesthesia will require a preoperative evaluation by an anesthesiologist to ensure they are medically fit for surgery and to discuss the anesthetic plan.
  \item \textbf{NPO Status}: Patients are instructed to be NPO (nil per os - nothing by mouth) for a specified period (typically 6-8 hours for solids, 2 hours for clear liquids) before surgery to prevent aspiration during anesthesia induction.
  \item \textbf{Informed Consent}: A detailed discussion of the procedure, its potential benefits, inherent risks, and alternative treatments is conducted, and the patient provides informed consent, acknowledging their understanding and agreement.
  \item \textbf{Prophylactic Antibiotics}: A single dose of prophylactic antibiotics (e.g., cefazolin) may be administered intravenously prior to incision to reduce the risk of surgical site infection.
\end{itemize}
Patients should also arrange for transportation home after the procedure, as they will not be able to drive due to the effects of anesthesia.

\section*{Contraindications \& Precautions}
Septoplasty, while generally safe, has specific contraindications and requires precautions in certain patient populations or clinical scenarios.

\textbf{Absolute Contraindications:}
\begin{itemize}
  \item Severe, uncontrolled systemic medical conditions (e.g., unstable angina, recent myocardial infarction, severe uncontrolled hypertension or diabetes, severe coagulopathy) that pose an unacceptable anesthetic or surgical risk.
  \item Active, severe nasal or sinus infection, which should be treated and resolved before surgery to minimize the risk of spreading infection and compromising healing.
  \item Uncorrectable severe coagulopathy or bleeding disorder that cannot be safely managed preoperatively.
  \item Patient refusal or inability to provide informed consent due to cognitive impairment without appropriate legal guardianship.
\end{itemize}

\textbf{Relative Contraindications and Precautions:}
\begin{itemize}
  \item \textbf{Young Age}: Septoplasty is generally deferred in children and adolescents whose nasal septum is still growing, typically until skeletal maturity (late teens), to avoid interfering with facial growth and potential recurrence of deviation. However, severe, symptomatic obstruction unresponsive to medical management may necessitate earlier, more conservative intervention.
  \item \textbf{Pregnancy}: Elective surgery is usually postponed until after delivery due to potential risks to the fetus from anesthesia and medications, especially during the first trimester.
  \item \textbf{Minor, Asymptomatic Deviation}: Surgery is not indicated for septal deviations that do not cause significant symptoms or functional impairment, as the risks outweigh the potential benefits.
  \item \textbf{Acute Rhinitis or Allergic Exacerbation}: Surgery may be postponed until these conditions are controlled to optimize mucosal health, reduce intraoperative bleeding, and improve postoperative healing.
  \item \textbf{Previous Nasal Surgery}: Prior surgery can complicate dissection due to scar tissue, requiring careful planning and potentially increasing the risk of mucosal tears or perforation.
  \item \textbf{Immunocompromised State}: Patients with compromised immune systems may be at higher risk for infection and impaired healing, requiring careful consideration and prophylactic measures.
\end{itemize}

\section*{Complications \& Risks}
While septoplasty is generally safe and effective, like any surgical procedure, it carries potential risks and complications that patients should be aware of.

\begin{itemize}
  \item \textbf{General Surgical Risks}:
    \begin{itemize}
      \item \textbf{Anesthesia Complications}: These can range from minor issues like nausea, vomiting, and sore throat to more serious, albeit rare, events such as allergic reactions, respiratory depression, or cardiovascular complications.
      \item \textbf{Bleeding (Hemorrhage)}: Postoperative epistaxis is the most common general complication, which can range from mild oozing to significant bleeding requiring intervention, such as repacking, cautery, or, rarely, reoperation.
      \item \textbf{Infection}: Although rare due to the vascularity of the nasal mucosa and prophylactic antibiotics, infection of the surgical site can occur, potentially leading to abscess formation (e.g., septal abscess) or cellulitis.
    \end{itemize}
  \item \textbf{Procedure-Specific Risks}:
    \begin{itemize}
      \item \textbf{Septal Hematoma}: A collection of blood between the septal cartilage and the mucoperichondrial flaps. If not drained promptly, it can compromise the blood supply to the cartilage, leading to cartilage necrosis and a saddle nose deformity.
      \item \textbf{Septal Perforation}: A hole in the septum, which can occur if both mucosal flaps are torn during surgery or if cartilage necrosis occurs postoperatively. Small perforations may be asymptomatic, while larger ones can cause crusting, whistling, recurrent epistaxis, or difficulty with nasal breathing.
      \item \textbf{Persistent or Recurrent Nasal Obstruction}: Incomplete correction of the deviation, scar tissue formation (synechiae), or recurrence of deviation due to cartilage memory can lead to ongoing breathing difficulties, potentially requiring revision surgery.
      \item \textbf{Change in Nasal Appearance (External Nasal Deformity)}: Although septoplasty is primarily an internal procedure, extensive cartilage removal or damage to the septal L-strut can rarely lead to a collapse of the nasal dorsum (saddle nose deformity) or changes in tip projection.
      \item \textbf{Adhesions (Synechiae)}: Scar tissue can form between the septum and the turbinates or lateral nasal wall, leading to new obstruction or difficulty with nasal hygiene.
      \item \textbf{Numbness}: Temporary or, rarely, permanent numbness of the upper front teeth, palate, or nasal tip due to nerve irritation or damage during flap elevation.
      \item \textbf{Cerebrospinal Fluid (CSF) Leak}: An extremely rare but serious complication if the cribriform plate is inadvertently breached during bony resection, requiring immediate repair to prevent meningitis.
      \item \textbf{Anosmia or Hyposmia}: Very rarely, a decrease or loss of the sense of smell can occur due to damage to the olfactory neuroepithelium.
      \item \textbf{Toxic Shock Syndrome}: An extremely rare but life-threatening complication associated with nasal packing, though less common with modern packing materials and shorter packing times.
    \end{itemize}
\end{itemize}

\section*{Postoperative Recovery Timeline}
Immediately after septoplasty, patients are transferred to the Post-Anesthesia Care Unit (PACU) for close monitoring as they recover from anesthesia. It is common to experience mild to moderate pain, nasal congestion, and some bloody drainage from the nose. Pain medication is administered as needed, typically oral analgesics. If nasal packing was used, it is usually soft and dissolvable or removed within 24-48 hours. Internal silicone splints, if placed, typically remain for 1-2 weeks. Patients are usually discharged home the same day, provided they are stable, have adequate pain control, and are accompanied by a responsible adult.

Over the first 1-2 weeks, patients will continue to experience significant nasal congestion due to internal swelling, crusting, and the presence of splints. Saline nasal rinses are highly recommended, often starting the day after surgery, to keep the nasal passages moist, aid in healing, and clear any crusting or dried blood. Patients are advised to avoid strenuous activities, heavy lifting, and nose blowing for at least one week to prevent bleeding and allow the septum to stabilize. Sneezing should be done with the mouth open to reduce intranasal pressure. Bruising around the eyes is uncommon but can occur, especially if combined with other procedures. Long-term recovery involves the complete resolution of internal swelling, which can take several weeks to months, leading to a gradual and progressive improvement in nasal breathing and overall comfort. Most patients can return to light activities within a few days and full activity (excluding contact sports) within 2-4 weeks.

\section*{Surgical Findings \& Pathology}
During septoplasty, the surgeon meticulously documents the intraoperative findings, which typically include the precise location and nature of the septal deviation (e.g., C-shaped deviation of the quadrangular cartilage to the right, with a bony spur along the maxillary crest impinging on the right inferior turbinate). The extent of cartilage and/or bone resected or reshaped, the integrity of the mucoperichondrial flaps, and the final alignment of the septum are also recorded. Any concomitant procedures, such as turbinate reduction, are noted.

For routine septoplasty, the resected septal cartilage and bone are generally not sent for pathological examination unless there is an unusual or suspicious finding, such as a mass, abnormal discoloration, or a history of prior malignancy. If tissue is sent, the pathology report typically confirms the presence of "benign nasal septal cartilage and/or bone." In rare instances where a suspicious lesion is encountered, the pathology report would provide a definitive diagnosis, guiding further management. The primary "pathology" in septoplasty is the anatomical deviation itself, which is corrected surgically rather than diagnosed histologically.

\section*{Expected Postoperative Course}
A normal and successful postoperative course following septoplasty is characterized by a predictable progression of healing and symptom resolution. Patients can expect mild to moderate pain for the first few days, which is typically well-managed with oral analgesics and gradually subsides. Nasal congestion will be prominent initially due to internal swelling, crusting, and the presence of any splints or packing. This congestion gradually improves over the first few weeks, with significant relief often noted after splint removal (if used) around 1-2 weeks post-surgery.

Serosanguinous drainage from the nose is common for the first few days. Patients will notice a gradual improvement in nasal breathing as the swelling resolves, with the full benefit of the surgery often becoming apparent within 1 to 3 months. The nasal mucosa will heal, and crusting will diminish with consistent saline rinses. There should be no signs of significant bleeding, infection, or external nasal deformity. The ultimate goal is a patent nasal airway, comfortable nasal breathing, and resolution of the preoperative symptoms, leading to an enhanced quality of life.

\section*{Postoperative Care \& Follow-up}
Postoperative care and diligent follow-up are essential to monitor healing, address any concerns, and ensure optimal long-term outcomes after septoplasty.

\begin{itemize}
  \item \textbf{Initial Postoperative Visit}: Typically scheduled 1-2 weeks after surgery for removal of any non-dissolvable sutures or internal silicone splints, if they were placed. The surgeon will also assess initial healing, remove any significant crusting, and provide further instructions on nasal hygiene.
  \item \textbf{Subsequent Follow-up Visits}: Additional visits may be scheduled at 1 month, 3 months, and sometimes 6 months post-surgery to monitor long-term healing, assess symptom improvement, and perform nasal endoscopy to check for any residual issues or complications like synechiae (adhesions) or recurrent deviation.
  \item \textbf{Nasal Endoscopy}: This is a routine part of follow-up to visualize the internal nasal structures, confirm the septum's position, and assess mucosal health and patency of the nasal passages.
  \item \textbf{Continued Nasal Hygiene}: Patients are usually advised to continue saline nasal rinses for several weeks to months to promote healing, prevent crusting, and maintain mucosal health. Topical nasal steroids may also be prescribed to reduce inflammation.
  \item \textbf{Activity Resumption}: Gradual return to normal activities is encouraged, with strenuous exercise and heavy lifting typically restricted for 2-4 weeks. Contact sports should be avoided for several months (e.g., 3-6 months) to protect the healing septum from trauma.
  \item \textbf{Warning Signs}: Patients are instructed to contact their surgeon immediately if they experience severe or worsening pain, heavy bright red bleeding that doesn't stop with pressure, fever (above 101.5$^{\circ}$F or 38.6$^{\circ}$C), foul-smelling nasal discharge, or any sudden changes in vision or mental status.
\end{itemize}
Adherence to these follow-up instructions is critical for a successful recovery and to achieve the best possible surgical outcome.

\section*{Epidemiology}
Nasal septal deviation is a highly prevalent anatomical condition, affecting a significant portion of the global population, though not all deviations are symptomatic. Studies suggest that anatomical septal deviation is present in approximately 80\% of adults, with varying degrees of severity. However, only a subset of these individuals experiences symptoms severe enough to warrant surgical intervention. Septoplasty remains one of the most common procedures performed by otolaryngologists worldwide, reflecting the high incidence of symptomatic nasal obstruction. Septal deviation can result from developmental factors, birth trauma, or, most commonly, external nasal trauma sustained during childhood or adulthood. There is no significant sex predilection, and it can affect individuals of all ages, though symptomatic presentation often becomes more pronounced in adulthood as nasal structures mature and other factors like allergic rhinitis or turbinate hypertrophy contribute to obstruction. The morbidity associated with untreated symptomatic septal deviation includes chronic nasal obstruction, sleep disturbances, recurrent sinusitis, and a significant reduction in quality of life. Mortality directly related to septal deviation or septoplasty is exceedingly rare, with most complications being manageable and non-life-threatening.

\section*{Outcomes \& Prognosis}
The outcomes of septoplasty are generally very favorable, with high rates of patient satisfaction and significant improvement in nasal breathing. Success rates, defined by subjective symptom improvement and objective findings of a straightened septum and patent airways, typically range from 85\% to 90\%. Functional outcomes, such as improved sleep quality, reduced snoring, decreased incidence of recurrent sinusitis, and enhanced exercise tolerance, are commonly reported. The prognosis for long-term relief from nasal obstruction is excellent for most patients, with sustained improvement over many years.

Factors influencing prognosis include the severity and complexity of the initial deviation, the presence of concomitant nasal pathologies (e.g., turbinate hypertrophy, allergic rhinitis, chronic rhinosinusitis), and the surgical technique employed. While recurrence of deviation is possible, especially in cases of severe cartilage memory or subsequent nasal trauma, it is relatively uncommon. Revision septoplasty may be necessary in a small percentage of patients (typically 5-10\%) who experience persistent or recurrent obstruction, often due to incomplete initial correction or scar tissue formation. Overall, septoplasty is a highly effective procedure for improving nasal function and significantly enhancing the patient's quality of life, as supported by numerous cohort studies and systematic reviews demonstrating its long-term efficacy.

\section*{References}
\begin{enumerate}
  \item MedlinePlus. Septoplasty. U.S. National Library of Medicine; 2024. Available from: \url{https://medlineplus.gov/septoplasty.html}
  \item Flint PW, Haughey BH, Lund VJ, et al., editors. Cummings Otolaryngology: Head \& Neck Surgery. 7th ed. Elsevier; 2021.
  \item American Academy of Otolaryngology--Head and Neck Surgery. Clinical Practice Guideline: Nasal Obstruction. Otolaryngology--Head and Neck Surgery. 2018;159(1\_suppl):S1-S33.
  \item Krouse JH, Christmas DA, Brown CL. Septoplasty: A Review of Indications, Techniques, and Outcomes. Otolaryngologic Clinics of North America. 2019;52(3):437-446.
\end{enumerate}

\clearpage